\documentclass[10pt]{amsart}
\usepackage[T1]{fontenc}
\usepackage{lmodern}
\usepackage[margin=1in]{geometry}
\usepackage{enumitem}
\setlist{itemsep=1pt,topsep=2pt,parsep=0pt}
\usepackage{amsmath,amssymb,array,longtable,booktabs}
\usepackage[hidelinks,pdftitle={Review: The Dirichlet tensor revision (uid452)},pdfauthor={Rick}]{hyperref}
\newcommand{\file}[1]{{\footnotesize\texttt{\detokenize{#1}}}}
\newcommand{\WIPHASH}{\texttt{848fc16}}
\newcommand{\tri}{\mathbin{\triangleleft}}
\newcommand{\Poly}{\mathbf{Poly}}
\newcommand{\Cat}{\mathbf{Cat}}
\newcommand{\Set}{\mathbf{Set}}
\newcommand{\Vect}{\mathbf{Vec}}
\newcommand{\Comon}{\mathbf{Comon}}
\newcommand{\E}{\mathcal{E}}
\newcommand{\ok}{\textsf{fixed}}

\title[Review: Dirichlet tensor revision (uid452)]{Review of the uid452 revision: \emph{The Dirichlet tensor of two polynomial comonoids is the product
internal category}}
\author{Rick}
\date{2026-10-24}

\begin{document}
\sloppy
\setlength{\parindent}{0pt}\setlength{\parskip}{3pt}
\maketitle
\vspace{-1em}
\begin{tabular}{@{}l>{\raggedright\arraybackslash}p{0.78\textwidth}@{}}
\textbf{Author} & Rick (grandparick20@gmail.com)\\
\textbf{Date} & 2026-10-24 (Wake 255)\\
\textbf{Recipient} & Scot MacBeth (scot.macbeth20@gmail.com); cc Robin Langer\\
\textbf{Title} & Review of the uid452 revision of the Dirichlet-tensor note\\
\textbf{WIP commit} & This PDF corresponds to WIP commit \WIPHASH\ (\texttt{grandpa-rick/work-in-progress},
\file{reviews/}), which holds this review's source. The hash was stamped in a follow-up commit, which adds this PDF.\\
\textbf{Reviewed} & \file{2026-10-22-day-convolution-topos-uniform.pdf}, content \texttt{32fe306}, stamp
\texttt{255a884}, 8 pp.\ (email UID 452). I compared it against the pre-452 version and against my uid440 report
(WIP \texttt{0237120}).\\
\textbf{Sources} & Shapiro--Spivak 2305.00167 (SS), full text re-read today for Def.\,3.7, 3.13, 4.1, Lem.\,4.3,
5.5, Thm.\,4.5, 5.6 and Cor.\,5.12. Lean file not run.\\
\end{tabular}

\medskip
\textbf{Bottom line.} Of the 13 uid440 items, 12 are fixed and 1 is partly fixed (M1: one false display
survives on page 1). Your question (1): \textbf{OK}. The equality is genuine, but three sentences should say why
(SHOULD S-A). Your question (2): \textbf{a lemma}, not a standalone note. The revision introduces \textbf{3 further
SHOULD} (S-B--S-D) and \textbf{3 NIT}, so 4 new SHOULD in all and no new MUST. One of them (S-B, Rem.\,3.1) comes
from a sentence of mine in uid440 that was too quick. The widened
hypothesis is enough, and the walking-arrow example is still consistent with Theorem 1.1.

\section{The uid440 items}

\begin{center}\small
\begin{longtable}{@{}l l p{0.70\textwidth}@{}}
\toprule
item & status & evidence in the revision\\ \midrule
M1 & \textbf{partly} & Thm.\,1.2 is gone. Example 4.1 is the walking-arrow witness, and the remark ``monoidal,
non-cartesian'' is in. SS is described with cofunctors in \S1 and \S4. \textbf{But} the \S1 display
$\Comon(\Poly,\tri,y)\simeq\Cat$ (p.\,1) survives. So do the chain ``containers $\simeq$ directed containers $\simeq$
polynomial comonoids $\simeq$ small categories'' (p.\,1) and ``$\simeq\cdots\simeq\Cat$'' (p.\,8). With comonoid
morphisms, the category is $\cong\Cat^\sharp$ (SS Cor.\,5.12), and \S4 of your own note shows that $\Cat^\sharp$ is
not $\Cat$. Fix: write $\cong\Cat^\sharp$, or say ``bijection on objects''. Separately, containers (polynomials) are
not equivalent to directed containers (comonoids). Drop the first $\simeq$.\\
M2 & \ok & $(\star)$ is now ``products commute with pullbacks''. GK is cited only for $\Poly(\E)$, $\tri$ over an
LCCC (p.\,4). The sharpening narrative and the $\Delta_A^*$ gloss are gone.\\
M3 & \ok & Retitled. Stated over any finitely complete $\E$ with exponentiable $\mathrm{dom}$. Def.\,2.2 makes no Day
claim over $\E$. \S7 and Table 1 are rewritten.\\
S1 & \ok & NS Prop.\,8.79 is cited by number, and the result is called a corollary of SS 4.5 + 5.6. The SS
Thm.\,4.5 eq.\,(16) and Ex.\,4.6 locators are right, and Ex.\,2.9 is gone. The Lem.\,5.5 alternative was added, but
see new item S-B.\\
S2 & \ok & The transpose of $\mathrm{cod}$ (\S2.2), $\iota$ on positions and naturality (\S2.3), the iso-class
remark (\S1), and the citation of the \emph{proof} of SS 5.6 (Step 3) are all in. The wording is still imprecise in
places; see \S2.\\
S3 & \ok & Lucyshyn-Wright is removed. The Vec paragraph is softened, but it now contains an inaccurate sentence
(new S-C).\\
S4 & \ok & The witness is framed as an illustration, both caveats are stated, ``first Lean certificate'' is dropped,
and the URL is given.\\
N1--N6 & \ok & $\Sigma_{A_p}\Pi_{s_p}B_p^*X$ is fixed. $C_{\mathrm{pt}}$ is now the positive control. The [5,
Prop.\,5.2] citation is replaced by SS 5.6/Rem.\,5.7. ``Ours'' is consistent. Table row 1 is fixed. Eq.\,(16) is not
reproduced.\\
\bottomrule
\end{longtable}
\end{center}

\section{Question (1): is Step 3's ``on the nose'' a genuine equality? Verdict: OK (with a SHOULD on wording)}

\textbf{What the objects and morphisms are.} SS Def.\,3.13 is explicit. The \emph{objects} of $\Poly_\E$ are
polynomials, i.e.\ actual exponentiable maps $s\colon P^*\to P$ of $\E$. They are not isomorphism classes.
\emph{Morphisms} are isomorphism classes of diagrams (SS Def.\,3.7), ``where the relevant isomorphisms are between
choices of the pullback $P\times_QQ^*$''. An object of $\Comon(\Poly_\E)$ is a triple $(c,\varepsilon,\delta)$. So
$c\otimes d=(C\times D)^\sharp$ in $\Comon$ means three things: (a) the two exponentiable maps are equal, (b) the
two counits are equal iso-classes, and (c) the two comultiplications are equal iso-classes.

\textbf{(a) holds literally, given one convention.} $c\otimes d$ is $\mathrm{dom}_C\times\mathrm{dom}_D$, formed
with the chosen binary products of $\E$ (SS Def.\,4.1). Suppose $C\times D$ means the \emph{componentwise} internal
category, $(C\times D)_i:=C_i\times D_i$, formed with the \emph{same} chosen products. Then its source map is
literally the same arrow of $\E$. Theorem 1.1 does specify the components, so this is fine. But ``the product
$C\times D$ in $\Cat(\E)$'' names an object only up to unique isomorphism. For any other product, the statement is
an isomorphism, not an equality. Say once: ``all limits are the chosen ones, and $C\times D$ is the componentwise
product formed with them''.

\textbf{(c) holds as an equality of iso-classes.} Both $\delta_{c\otimes d}=\iota\circ(\delta_c\otimes\delta_d)$
and $\delta_{(C\times D)^\sharp}$ are morphisms $c\otimes d\to(c\otimes d)\tri(c\otimes d)$, between literally the
same two objects, by (a).
\begin{itemize}[leftmargin=*]
\item \emph{Positions.} Both forward maps land in the one chosen object $\Sigma\Pi_{\mathrm{dom}}(\cdots)$. Each is
  the transpose of a map out of $C_1\times D_1$. Transposes are unique, so the two are equal as arrows of $\E$
  exactly when $\mathrm{cod}_{c\otimes d}=\mathrm{cod}_C\times\mathrm{cod}_D$. That is your naturality line. No
  isomorphism is involved.
\item \emph{Directions.} The backward leg of $\iota$ is the canonical comparison
  $\kappa\colon(C_1\times D_1)\times_{C_0\times D_0}(C_1\times D_1)\to(C_1\times_{C_0}C_1)\times(D_1\times_{D_0}D_1)$,
  modulo the choice of pullback. That choice is exactly what Def.\,3.7 quotients by. The backward leg of
  $\delta_c\otimes\delta_d$ is $\mathrm{comp}_C\times\mathrm{comp}_D$. So the backward leg of $\delta_{c\otimes d}$
  is $(\mathrm{comp}_C\times\mathrm{comp}_D)\circ\kappa$. And the composition of the componentwise product category
  is \emph{defined} on the chosen pullback of products as $(\mathrm{comp}_C\times\mathrm{comp}_D)\circ\kappa$.
  There is nothing else it could be. So the two backward legs are the same arrow, not isomorphic arrows.
\end{itemize}
The SS dictionary (proof of Thm.\,5.6) reads $\delta^\sharp$ through an isomorphism
$\coprod\coprod c[\delta_1(x)(y)]\cong C^*\,{}_t\!\times_c C^*$. That too is a change of pullback representative,
absorbed by Def.\,3.7.

\textbf{(b) holds, but the one visible isomorphism is in your definition of $\varepsilon$.} $\varepsilon_c\otimes
\varepsilon_d$ lands in $y\otimes y=(1\times1\to1\times1)$. This equals $y=(1\to1)$ only if the chosen $1\times1$ is
$1$. Your \S3 writes ``into $y\otimes y\cong y$''. So $\varepsilon_{c\otimes d}$ is $\lambda\circ(\varepsilon_c
\otimes\varepsilon_d)$, with $\lambda$ the unitor. This is the correct definition of the tensor of comonoids in a
normal duoidal category. Its iso-class has forward leg $!$ and backward leg $\mathrm{id}_C\times\mathrm{id}_D$, which
is literally $\varepsilon_{(C\times D)^\sharp}$. So (b) is an equality too.

\textbf{Verdict on (1): OK.} The equality is genuine. No coherence isomorphism is suppressed, because the only two
candidates, $\kappa$ and $\lambda$, are respectively part of the \emph{definition} of $\mathrm{comp}_{C\times D}$
and part of the \emph{definition} of $\varepsilon_{c\otimes d}$. But the text does not yet say this:
\begin{enumerate}[label=\textbf{S-A.},leftmargin=*]
\item \textbf{(SHOULD) Three wording fixes.}
  (i) In \S1, ``an equality of polynomials is a genuine equality of such classes'' conflates objects with morphisms.
  Polynomials are objects, and equality of objects is literal (given chosen products). Only equality of the structure
  \emph{morphisms} is an equality of iso-classes.
  (ii) In Step 3 (Composition), ``$\mathrm{comp}_{c\otimes d}=\mathrm{comp}_C\times\mathrm{comp}_D$ under the
  canonical isomorphism $(\star)$'' reads exactly like a suppressed coherence iso, which is what you asked me to
  check. Replace it with: ``$\kappa$ is the backward leg of $\iota$. By definition of the componentwise product,
  $\mathrm{comp}_{C\times D}=(\mathrm{comp}_C\times\mathrm{comp}_D)\circ\kappa$. Hence the backward legs are equal.''
  (iii) Add the chosen-limits convention of (a). Mention $\lambda$ in the definition of $\varepsilon_{c\otimes d}$
  (or choose $1\times1=1$).
\end{enumerate}
A small simplification follows from (a)--(c). Once the triple equals $(C\times D)^\sharp$, the comonoid laws are
immediate: $(C\times D)^\sharp$ is a comonoid by the forward direction of SS Thm.\,5.6, applied to the internal
category $C\times D$. The appeal to the one-by-one matching inside SS's proof is now unnecessary. (I asked for it in
uid440 when Rem.\,3.1 argued differently.)

\section{Question (2): standalone note or lemma? Verdict: a lemma}

Blunt version. The Set case is NS Prop.\,8.79, word for word. The internal case unwinds SS Def.\,4.1, Thm.\,4.5 and
Thm.\,5.6, and its proof is half a page (Steps 1--3, which \S2 above compresses further). Example 4.1 makes concrete
a sentence that NS already print. Table 1 and \S7 are accounting. None of this is wrong now, and the accounting is
honest. But 8 pages of framing around a half-page corollary will read to a referee as a remark. A referee will also
note that SS's Lem.\,5.5 technique makes the general-$\E$ step routine. The rhetoric in \S1 (``a question that
inverted into its own answer'') and \S7 (``what this buys the chain'') promises more than the result delivers.

Recommendation: state it as a lemma, with Example 4.1 as the following remark, in whatever unit uses internal
polynomial comonoids in presheaf or sheaf models (your Graph models are the natural host). It would become a
note-grade standalone only with a second result of its own. Two candidates, both unchecked for novelty:
(a) describe the actual categorical product in $\Cat^\sharp(\E)$, or show it fails to exist over general $\E$. NS
remark that $\mathbf2\times^\sharp\mathbf2$ has infinitely many objects. An internal version very likely needs infinite
colimits in $\E$, so the base hypothesis would finally do real work.
(b) Prove the Day-convolution identification over Grothendieck $\E$ in a stated sense (uid440 option (b)).

\section{New issues introduced by the revision}

\begin{enumerate}[label=\textbf{S-\Alph*.},leftmargin=*,start=2]
\item \textbf{(SHOULD) Rem.\,3.1, the ``three-line alternative'', has a gap. It contradicts the note's own \S1 and
  \S6.} ``Pointwise, the claim is the Set case [2, Prop.\,8.79], applied at each object'' is unjustified. SS
  Lem.\,5.5 transports structures to $\widehat{\mathcal A}$. But there $c\tri c$ is \emph{not} computed at each
  stage: SS Ex.\,3.21 and diagram (29) use $\mathrm{Hom}_{\widehat{\mathcal A}}(c[x],C)$, not
  $\mathrm{Hom}_\Set(c_a[x],C_a)$. Evaluation at $a$ does not preserve $\Pi_{\mathrm{dom}}$. Your \S1 (``not computed
  objectwise in a presheaf topos'') and \S6 caveat say the same. So the Set case at each $a$ concerns a different
  comonoid, and the reduction proves nothing beyond ``it suffices to prove Theorem 1.1 in $\widehat{\mathcal A}$''.
  In $\widehat{\mathcal A}$ you would redo Step 3. This is my fault: uid440 S1 said ``in presheaves the category side
  is pointwise'', which is true but does not finish the argument. Delete Rem.\,3.1, or cut it to that one true
  sentence.
\item \textbf{(SHOULD) The \S1 Vec paragraph: ``$\Poly(\Vect)$ has no standard construction'' is not right, and the
  precise fact is better.} $\Vect$ is finitely complete, so SS applies verbatim. But in $\Vect$ \emph{the
  exponentiable maps are exactly the isomorphisms}. Proof: $s^*\colon\Vect/A\to\Vect/B$ must preserve the coproduct
  $A\oplus A\xrightarrow{+}A$ of $\mathrm{id}_A$ with itself. The comparison $B\oplus B\to(A\oplus
  A)\times_AB$, $(b,b')\mapsto(sb,sb',b+b')$, sends $(b,-b)$ to $(sb,-sb,0)$, so injectivity forces $s$ to be
  injective. Surjectivity needs every $(a,a',b)$ with $a+a'=sb$ to have $a\in\operatorname{im}s$, and $a$ is
  arbitrary, so $s$ is surjective. Hence the internal categories with exponentiable $\mathrm{dom}$ are the discrete
  ones (2-term complexes with zero arrow part). Theorem 1.1 holds over $\Vect$ but is degenerate there. That is the
  exact location of the ``Vec obstruction'': it is a shortage of exponentiable maps, not a lack of a construction.
\item \textbf{(SHOULD) Example 4.1 describes the projections as functors.} ``$\pi_C,\pi_D$ are the cofunctors
  $(f,h)\mapsto f$ and $(f,h)\mapsto h$ on morphisms'' goes the wrong way. A cofunctor acts \emph{backward} on
  morphisms. $\pi_C=\lambda\circ(\mathrm{id}_c\otimes\varepsilon_d)$ lifts $f\colon x\to x'$ at $(x,y)$ to
  $(f,\mathrm{id}_y)$, because the backward leg of $\varepsilon_d$ is $\mathrm{id}$. That lift is the one the
  argument actually uses two lines later, when $H$ must lift $(a,\mathrm{id})$. Rewrite it as ``lifting $f$ at
  $(x,y)$ to $(f,\mathrm{id}_y)$''.
\end{enumerate}
\textbf{NITs.} (N-a) The Lean URL in \S6 runs off the right margin and is visibly cut (``\dots
convolution-grap''). Use \verb|\url| with \texttt{xurl}, or break it. (N-b) Example 4.1 uses italic $E$ for the test
category next to the calligraphic base $\E$. Rename it $T$. (N-c) Step 2 says the content of $\iota$ on directions
``is'' $(\star)$. Say ``is the comparison map $\kappa$ of $(\star)$''. That matters for S-A(ii).

\textbf{Consistency checks (no issue found).}
\begin{itemize}[leftmargin=*]
\item \emph{Is ``any finitely complete $\E$'' enough?} Yes. $\mathrm{dom}_C\times\mathrm{dom}_D=(\mathrm{dom}_C\times
  1)\circ(1\times\mathrm{dom}_D)$ is a composite of pullbacks of exponentiable maps, hence exponentiable (SS
  Lem.\,4.3). $\tri$ is monoidal on $\Poly_\E$ for finitely complete $\E$ (SS Thm.\,3.26). $\iota$ and normal
  duoidality exist there (SS Thm.\,4.5). The comonoid--category bijection holds there (SS Thm.\,5.6). Steps 1--3 use
  only chosen finite limits and the universal property of $\Pi_{\mathrm{dom}}$, which exists by hypothesis. Nothing
  needs local cartesian closure, coproducts or a subobject classifier.
\item \emph{Walking arrow against Theorem 1.1.} The theorem now asserts only an identification of the comonoid
  $c\otimes d$, which is an object of $\Comon$. It asserts no universal property. Example 4.1 refutes only
  universality among cofunctors. The abstract, \S1, \S4 and \S7 all say ``monoidal, non-cartesian''. I re-ran the
  witness by hand. $H(e)=(0,0)$ and $H(x)=H(z)=(1,1)$ are forced. $\pi_CH=F$ forces the lift of $(a,\mathrm{id})$ at
  $e$ to be $u$. But $H(\mathrm{cod}\,u)=(1,1)\neq(1,0)$. Contradiction, as stated. The surviving false sentence is
  the \S1 display in M1, not the theorem.
\item \emph{Cat$(\E)$ product.} The sketch argument in \S4 is correct. The projections are internal functors.
\end{itemize}

\end{document}
